\documentclass[aspectratio=169]{beamer}

\usetheme{metropolis}
\usepackage[utf8]{inputenc}
\usepackage[T1]{fontenc}
\usepackage[spanish]{babel}
\usepackage{xcolor}
\usepackage{booktabs}
\usepackage{array}
\usepackage{tabularx}
\usepackage{tikz}

% Paleta personalizada
\definecolor{navy}{HTML}{1A2744}
\definecolor{teal}{HTML}{028090}
\definecolor{accent}{HTML}{C74B2A}
\definecolor{gold}{HTML}{D4A843}
\definecolor{verde}{HTML}{2D6A4F}
\definecolor{gris}{HTML}{6C757D}

% Aplicar paleta al tema metropolis
\setbeamercolor{normal text}{fg=navy, bg=white}
\setbeamercolor{alerted text}{fg=accent}
\setbeamercolor{example text}{fg=teal}
\setbeamercolor{structure}{fg=navy}
\setbeamercolor{frametitle}{bg=navy, fg=white}
\setbeamercolor{title separator}{fg=teal}
\setbeamercolor{progress bar}{fg=teal}
\setbeamercolor{progress bar in head/foot}{fg=teal}
\setbeamercolor{progress bar in section page}{fg=teal}

% Bloques personalizados
\setbeamercolor{block title}{bg=teal, fg=white}
\setbeamercolor{block body}{bg=teal!10, fg=navy}

\setbeamercolor{block title alerted}{bg=accent, fg=white}
\setbeamercolor{block body alerted}{bg=accent!10, fg=navy}

\setbeamercolor{block title example}{bg=verde, fg=white}
\setbeamercolor{block body example}{bg=verde!10, fg=navy}

% Bloque oro para facciones condicionales
\newenvironment{goldblock}[1]{%
  \setbeamercolor{block title}{bg=gold, fg=navy}
  \setbeamercolor{block body}{bg=gold!15, fg=navy}
  \begin{block}{#1}
}{%
  \end{block}
}

\newenvironment{grayblock}[1]{%
  \setbeamercolor{block title}{bg=gris, fg=white}
  \setbeamercolor{block body}{bg=gris!15, fg=navy}
  \begin{block}{#1}
}{%
  \end{block}
}

\title{Simulación: Sesión de Concejo Municipal}
\subtitle{Proyecto de Acuerdo: Actualización Catastral en los municipios de Cundinamarca}
\author{Prof.~Andrés Pérez Coronado, PhD}
\institute{Especialización en Gobernanza y Desarrollo Territorial\\Universidad del Rosario}
\date{Abril 2026}

\begin{document}

% ----------------------------------------------------------------------
% SLIDE 1 — PORTADA
% ----------------------------------------------------------------------
\maketitle

% ----------------------------------------------------------------------
% SLIDE 2 — OBJETIVO DEL EJERCICIO
% ----------------------------------------------------------------------
\begin{frame}{Objetivo del ejercicio}
  \small
  Simular una sesión plenaria de concejo municipal donde ustedes, como
  concejales de distintos municipios de Cundinamarca, deben \textbf{debatir y
  votar} un proyecto de acuerdo sobre la actualización catastral
  multipropósito.
  \vspace{0.6em}

  \begin{enumerate}
    \item \textbf{\textcolor{teal}{Argumentar con datos}} --- construir
    posiciones basadas en evidencia, no en intuición.

    \item \textbf{\textcolor{accent}{Negociar intereses contrapuestos}}
    --- encontrar soluciones que equilibren recaudo fiscal y bienestar
    ciudadano.

    \item \textbf{\textcolor{verde}{Pensar como gerentes públicos}} ---
    aplicar herramientas de gerencia estratégica a un problema real.
  \end{enumerate}
\end{frame}

% ----------------------------------------------------------------------
% SLIDE 3 — EL PROBLEMA (CONTEXTO REAL)
% ----------------------------------------------------------------------
\begin{frame}{El problema: actualización catastral en Cundinamarca}
  {\color{accent}\textbf{Contexto real --- esto está pasando ahora}}
  \vspace{0.4em}

  \small
  \begin{itemize}\itemsep2pt
    \item La Agencia Catastral de Cundinamarca actualizó los avalúos en
    \textbf{76 municipios}.
    \item Se han recibido más de \textbf{12.000 reclamaciones} de
    ciudadanos inconformes.
    \item Algunos predios rurales vieron incrementos de hasta \textbf{300\%}
    en su avalúo catastral.
    \item El impuesto predial se calcula sobre el avalúo: si el avalúo
    sube, el predial sube.
    \item Campesinos y pequeños propietarios alertan que no pueden pagar.
    \item Los municipios necesitan el recaudo para financiar inversión
    social.
    \item La Gobernación activó la \textit{Ruta del Diálogo Catastral}
    para atender reclamos.
    \item \textbf{13 zonas} fueron declaradas críticas con ventanillas
    exclusivas.
  \end{itemize}
  \vspace{0.3em}
  {\scriptsize \textit{Fuente: Gobernación de Cundinamarca, Ruta del
  Diálogo Catastral, 2024--2026.}}
\end{frame}

% ----------------------------------------------------------------------
% SLIDE 4 — PROYECTO DE ACUERDO (FICTICIO)
% ----------------------------------------------------------------------
\begin{frame}{Proyecto de Acuerdo No.~001 de 2026}
  \scriptsize
  \begin{block}{Articulado}
    \begin{description}\itemsep3pt
      \item[\textcolor{navy}{Art.~1}] Adóptese la actualización catastral
      multipropósito como base para el cálculo del impuesto predial a
      partir del año gravable \textbf{2027}.

      \item[\textcolor{navy}{Art.~2}] Establézcase un \textbf{régimen de
      transición de 3 años} con incrementos graduales: Año~1 (30\%),
      Año~2 (60\%), Año~3 (100\%) del nuevo avalúo.

      \item[\textcolor{navy}{Art.~3}] Créese un \textbf{fondo de alivio
      tributario} para predios rurales menores a 5~ha cuyos propietarios
      estén en SISBEN A o B.

      \item[\textcolor{navy}{Art.~4}] Destínese el \textbf{40\% del mayor
      recaudo} a inversión en vías terciarias, agua potable y
      conectividad rural del municipio.

      \item[\textcolor{navy}{Art.~5}] La administración municipal rendirá
      \textbf{informe trimestral} al Concejo sobre recaudo, reclamos y
      ejecución del fondo de alivio.
    \end{description}
  \end{block}
  \vspace{0.3em}
  {\color{accent}\textbf{\scriptsize ⚠ DOCUMENTO FICTICIO PARA EJERCICIO
  ACADÉMICO}}
\end{frame}

% ----------------------------------------------------------------------
% SLIDE 5 — LOS ROLES (4 FACCIONES)
% ----------------------------------------------------------------------
\begin{frame}{Los roles: 4 facciones en el Concejo}
  \tiny
  \begin{columns}[t]
    \begin{column}{0.49\textwidth}
      \begin{exampleblock}{A FAVOR (verde · 35\%)}
        \textit{«Apoyar la actualización. El municipio necesita recursos
        y el predial está desactualizado hace 10+ años.»}
        \vspace{0.3em}\\
        \textbf{Argumentos:} rezago catastral, equidad tributaria,
        financiación del Plan de Desarrollo.
      \end{exampleblock}

      \begin{goldblock}{A FAVOR CON CONDICIONES (gold · 25\%)}
        \textit{«Aprobar solo si se modifican artículos: ampliar
        transición a 5 años, más beneficiarios en el fondo de alivio,
        auditoría externa.»}
        \vspace{0.3em}\\
        \textbf{Argumentos:} gradualidad, protección social, control y
        transparencia.
      \end{goldblock}
    \end{column}

    \begin{column}{0.49\textwidth}
      \begin{alertblock}{EN CONTRA (rojo · 30\%)}
        \textit{«Rechazar. Los avalúos son excesivos, no consultan la
        capacidad de pago.»}
        \vspace{0.3em}\\
        \textbf{Argumentos:} impacto social, errores en avalúos,
        capacidad de pago, riesgo electoral.
      \end{alertblock}

      \begin{grayblock}{INDECISOS (gris · 10\%)}
        \textit{«No tienen posición definida. Su voto define el
        resultado. Representan municipios pequeños con poca
        información.»}
        \vspace{0.3em}\\
        \textbf{Argumentos:} necesitan datos, quieren entender el
        impacto real.
      \end{grayblock}
    \end{column}
  \end{columns}
\end{frame}

% ----------------------------------------------------------------------
% SLIDE 6 — DATOS PARA EL DEBATE
% ----------------------------------------------------------------------
\begin{frame}{Datos para argumentar}
  \scriptsize
  \begin{tabularx}{\textwidth}{@{}X l l@{}}
    \toprule
    \textbf{Indicador} & \textbf{Dato} & \textbf{Fuente} \\
    \midrule
    Municipios actualizados en Cundinamarca & 76 de 116 & Gobernación, 2026 \\
    Reclamaciones ciudadanas recibidas & Más de 12.000 & Agencia Catastral \\
    Zonas declaradas críticas & 13 zonas & Ruta del Diálogo \\
    Incremento promedio en avalúos rurales & Entre 80\% y 300\% & Reportes de prensa \\
    Rezago catastral promedio antes de actualización & 10--15 años & IGAC / DNP \\
    \% del predial en ingresos propios municipales & 30--60\% según municipio & DNP \\
    Hogares rurales Cund.\ sin internet & Solo 41,9\% conectados & DANE, ENTIC 2024 \\
    \bottomrule
  \end{tabularx}
  \vspace{0.6em}

  \begin{block}{\small Instrucción}
    \scriptsize Cada facción debe usar \textbf{al menos 2 datos} de esta
    tabla en sus argumentos.
  \end{block}
\end{frame}

% ----------------------------------------------------------------------
% SLIDE 7 — REGLAS DE LA SIMULACIÓN
% ----------------------------------------------------------------------
\begin{frame}{Reglas de la simulación}
  \small
  \begin{enumerate}\itemsep3pt
    \item \textbf{\textcolor{teal}{Formación de facciones (5 min)}}
    --- se asignan roles; cada grupo elige vocero y prepara posición.

    \item \textbf{\textcolor{teal}{Presentación de posiciones (3 min
    c/u)}} --- cada facción expone; deben usar datos.

    \item \textbf{\textcolor{teal}{Debate abierto (15 min)}} ---
    interpelaciones entre facciones; se permite proponer modificaciones.

    \item \textbf{\textcolor{teal}{Negociación (10 min)}} --- las
    facciones negocian cambios para lograr votos.

    \item \textbf{\textcolor{teal}{Votación nominal (5 min)}} --- cada
    concejal vota SÍ, NO o ABSTENCIÓN. \textbf{Mayoría simple}.

    \item \textbf{\textcolor{teal}{Debrief (10 min)}} --- reflexión
    grupal.
  \end{enumerate}
\end{frame}

% ----------------------------------------------------------------------
% SLIDE 8 — PREGUNTAS GUÍA
% ----------------------------------------------------------------------
\begin{frame}{Preguntas guía para preparar su posición}
  \scriptsize
  \begin{enumerate}\itemsep2pt
    \item ¿Cuánto recauda su municipio por predial hoy y cuánto necesita
    para el Plan de Desarrollo?
    \item ¿Qué porcentaje de los predios de su municipio son rurales
    vs.\ urbanos?
    \item ¿Cuántos ciudadanos han reclamado en su municipio?
    \item ¿El régimen de transición de 3 años es suficiente o debería
    ser más largo?
    \item ¿El fondo de alivio tributario cubre a quienes realmente lo
    necesitan?
    \item ¿Qué pasa si se rechaza el proyecto? ¿De dónde sale la plata
    para inversión?
    \item ¿Qué modificaciones propondría usted a los artículos del
    proyecto?
    \item ¿Cómo le explicaría su voto a los ciudadanos de su
    municipio?
  \end{enumerate}
\end{frame}

% ----------------------------------------------------------------------
% SLIDE 9 — ¡COMIENZA LA SESIÓN!
% ----------------------------------------------------------------------
{
\setbeamercolor{background canvas}{bg=navy}
\setbeamercolor{normal text}{fg=white}
\begin{frame}[plain]
  \centering
  \vspace{0.6cm}
  {\color{white}\Huge\textbf{¡Comienza la sesión!}}
  \vspace{0.6cm}

  {\color{gold}\Large Paso 1: Facciones --- reúnanse y preparen su
  posición}
  \vspace{1.4cm}

  {\color{white}\Huge\textbf{5 minutos}}
  \vspace{1.2cm}

  {\color{white}\normalsize
  Elijan vocero $\cdot$ Definan 3 argumentos $\cdot$ Usen los datos de
  la slide 6}
\end{frame}
}

% ----------------------------------------------------------------------
% SLIDE 10 — REFLEXIÓN / DEBRIEF
% ----------------------------------------------------------------------
\begin{frame}{Reflexión: ¿Qué aprendimos?}
  \tiny
  \begin{columns}[t]
    \begin{column}{0.49\textwidth}
      \begin{block}{Sobre el uso de datos}
        ¿Fue más fácil argumentar cuando tenían cifras concretas? ¿Qué
        datos les faltaron?
      \end{block}
      \begin{exampleblock}{Sobre la gerencia estratégica}
        ¿El proyecto resuelve un problema o crea otros? ¿Cómo lo
        mejorarían?
      \end{exampleblock}
    \end{column}

    \begin{column}{0.49\textwidth}
      \begin{alertblock}{Sobre la negociación política}
        ¿Cómo lograron (o no) convencer a los indecisos? ¿Qué
        concesiones hicieron?
      \end{alertblock}
      \begin{goldblock}{Sobre su rol como concejales}
        ¿Toman decisiones así en sus concejos reales? ¿Tienen acceso a
        datos para debatir?
      \end{goldblock}
    \end{column}
  \end{columns}
  \vspace{0.8em}
  \centering
  \textit{\color{navy}\small «La IA y los datos son herramientas. Las
  decisiones las toman ustedes.»}
\end{frame}

\end{document}
